\documentclass[book.tex]{subfiles}

\begin{document}

\chapter{Physika Primer}

\label{chap:physika}
\noindent\rule{11cm}{0.4pt} \\
\textbf{Prerequisites:}  \autoref{chap:math_basics}   \\
\textbf{Difficulty Level:} * \\
\noindent\rule{11cm}{0.4pt}
\vspace{5mm}


Throughout this book, we use pseudocode to represent various differentiable programs of interest. We code our pseudocode in a language called Physika (greek for physics). Physika is a fully differentiable, dependently typed language. This combination of features allows Physika to model physical theories succinctly and effectively. In this chapter we provide a brief introduction to Physika. 

\section{An Introduction to Physika by Example}
Much of Physika's syntax follows that of Python. For example, here's how you define a function that adds two numbers.

\begin{lstlisting}[language=python]
# This is a comment.
def add(x: $\mathbb{R}$, y: $\mathbb{R}$) $\to$ $\mathbb{R}$:
  x + y
\end{lstlisting}

Notice that we don't define an explicit return type. Physika follows the convention that the last expression in a function is its return value. This allows for simple mathematically crisp definitions.

Physika is a typed language with a rich and descriptive type system for annotating function inputs and outputs. We allow variables to have type $\mathbb{R}$. We can't directly represent infinite decimal place real numbers on an actual computer. But we can symbolically manipulate values in $\mathbb{R}$. We can also specialize this function to operate on a subset of $\mathbb{R}$ like the 32 bit floating point numbers $\textrm{Float32}$. In practice, it will be clear how to specialize a function that operates on the real numbers to a computationally realizable version, but if there are subtleties we will mention them with examples.

We define the type of an array as
\begin{lstlisting}[language=python]
x : $\mathbb{R}$[n]
\end{lstlisting}
Multidimensional arrays are defined similarly
\begin{lstlisting}[language=python]
y : $\mathbb{R}$[n, m, o]
\end{lstlisting}
The type of a tensor is similar to the type for a multidimensional array, with a crucial difference that covariant and contravariant indices must be annotated in the type.
\begin{lstlisting}[language=python]
z : $\mathbb{R}$[+n, +m, -o]
\end{lstlisting}
In mathematical tensor notation we would write this as something like $z^{\alpha \beta}_{\gamma}$. Note that tensor $z$ is not the same as a multidimensional array since it depends on a choice of basis. It is possible to \textit{cast} a tensor object into a multidimensional array. The default cast will use the current basis of the tensor object.

Physika is dependently typed. What does that mean? To start, it means that the following function is valid.

\begin{lstlisting}[language=python]
def append(x: a, xs: a[n]) $\to$ a[n+1]:
  xs + [x]
\end{lstlisting}

Note that the type \inline{a} is a free variable. This is typically referred to as parameteric polymorphism in the programming language literature. You can think about \inline{a} as a placeholder that can be replaced with any other type like $\mathbb{R}$ or \inline{Float32}. Another interesting property of this definition is that the return type can depend on the provided arguments.

Let's now see how we can define a variable with a value

\begin{lstlisting}[language=python]
x : $\mathbb{R}$ = 5.12
\end{lstlisting}

Next, we consider a more complex example of matrix multiplication. Here is how we would code a matrix multiplication algorithm in Physika

\begin{lstlisting}[language=python]
def matmul(A: $\mathbb{R}$[n, m], B: $\mathbb{R}$[m, o]) $\to$ $\mathbb{R}$[n, o]:
  C: $\mathbb{R}$[n, o] = 0
  for i j k:
    C[i, j] += A[i,k]*B[k, j]
  C
\end{lstlisting}

We use the type system to infer that \lstinline{i} and \lstinline{j} above must respectively take values in \lstinline{[0,...,n-1]} and \lstinline{[0,...,o-1]}. Type inference allows us to make compact readable code for numerical applications. We can use type inference to simplify operations like Einstein summation. For example, an expression like
\begin{align}
    T^{\alpha}_\beta v^{\beta} 
\end{align}
which represents a matrix vector multiplication can be represented as a code snippet
\begin{lstlisting}[language=python]
T: $\mathbb{R}$[$+n$, $-m$]
v: $\mathbb{R}$[$+m$]
sum(for $\beta$ T[$\alpha$, $\beta$] * v[$\beta$])
\end{lstlisting}

We can also rely on Physika to do type inference on tensor shapes and covariances properly.
\begin{lstlisting}[language=python]
T: $\mathbb{R}$[+m, +n, -o]
U: $\mathbb{R}$[+o]
T[+$\alpha$, +$\beta$, -$\gamma$]U[$+\gamma$]: $\mathbb{R}$[$+\alpha$, $+\beta$]
\end{lstlisting}


Physika is a differentiable programming language, so differentiation operators are first class in the language. Phyika has two differentiation operators for forward and backwards mode differentiation
\begin{lstlisting}
T # Tangent Operator for forward mode
$\nabla$ # Gradient Operator for backwards mode
$\frac{d}{dt}$ # Alternative syntax for backwards mode
\end{lstlisting}

We now explore a slightly more complicated example using Physika to encode some simple motion of particles. We can encode 

\begin{lstlisting}[language=python]
x: $\mathbb{R}$ $\to$ $\mathbb{R}^3$
x(t) = (t, t, t)
v = $\frac{dx}{dt}$
\end{lstlisting}
What is the type of $v$? We write
\begin{lstlisting}[language=python]
v: $\mathbb{R}$ $\to$ $T \mathbb{R}^3$
\end{lstlisting}
Let's write down a type for the differentiation operator in this particular case.

\begin{lstlisting}[language=python]
$\frac{d}{dt}$: ($\mathbb{R}$ $\to$ $\mathbb{R}^3$) $\to$ $T \mathbb{R}^3$
\end{lstlisting}

Physika is capable of expressing more complex function spaces. Notably, we will make use of $L^2$ function space types for representing wave functions.
\begin{lstlisting}[language=python]
$\psi$: $L^2[\mathbb{R}^3, \mathbb{C}]$
\end{lstlisting}

The study of quantum mechanics encompasses not only the study of wave functions but of operators that act on them. For example, consider the momentum operator that acts on wave functions. This operator has type

\begin{lstlisting}[language=python]
$\hat{p}$: $L^2[\mathbb{R}^3]$ $\to$ $L^2[\mathbb{R}^3]$
\end{lstlisting}
We could try to write this type as
\begin{lstlisting}[language=python]
$\hat{p}$: $\mathcal{B}(L^2[\mathbb{R}^3])$
\end{lstlisting}


This model is wrong, since $\hat{p}$ is not necessarily bounded as an operator, but serves to show how Physika can represent the type of bounded operators on Hilbert space. We can represent common quantum mechanical operations such as bras, kets, and brakets with suitable types.
\begin{lstlisting}[language=python]
$|\psi(x)\rangle$ $\equiv$ $\psi: L^2[\mathbb{R}^3, \mathbb{C}]$
\end{lstlisting}

The type of a quantum mechanical bra is given by
\begin{lstlisting}[language=python]
$\langle \psi(x)|$ $\equiv$ $L^2[\mathbb{R}^3, \mathbb{C}] \to \mathbb{C}$
\end{lstlisting}

%\textcolor{red}{TODO: Introduce the notion of a \textit{type integral}, that is an integral over all values of a type.}

We can represent yet more complicated spaces. For example, to compute with path integrals, we need to represent the type of all evolution paths for a given quantum mechanical system.

\begin{lstlisting}[language=python]
p: [0, 1] $\to$ $L^2(\mathbb{R}^3)$
\end{lstlisting}


%\textcolor{red}{TODO: Add type of the path integral}

%\begin{lstlisting}[language=python]
%-i $\hbar$ 
%\end{lstlisting}

\section{Symbolic Manipulation}

For describing certain algorithms, it will be useful to perform symbolic manipulations.

\begin{lstlisting}[language=python]
x: Symbol
x$^4$ - 3x$^2$ + 13
\end{lstlisting}

Symbolic variables can be used to express and manipulate formulas.

\section{Parameterized Functions}

Layers in deep networks are typically represented as parameterized functions $f_\theta$. A parameterized function can be viewed as a family of concrete functions
\begin{align}
    \{ f_{\theta} | \theta \in \Theta \}
\end{align}

We write a parameterized layer function in Physika by using the \lstinline{class} keyword. 
\begin{lstlisting}[language=python]
class Layer(param$_1$: T$_1$,$\dotsc$):
  def $\lambda$(arg$_1$: A$_1$,$\dotsc$): R_1
    ...
\end{lstlisting}

We use the notation $\lambda$ to denote the function call to the parameterized function. The parameters of the function are set in the constructor. We can define a simple weighted linear layer using a parameterized function as follows

\begin{lstlisting}[language=python]
# Weighted Addition
class F($w_1$: $\mathbb{R}$, $w_2$: $\mathbb{R}$):
  def $\lambda$(x: $\mathbb{R}$, y: $\mathbb{R}$): $\mathbb{R}$:
    $w_1$ * x + $w_2$ * y
\end{lstlisting}

We can express the type of a parameterized function in a type signature as 

\begin{align}
    (\textrm{Input Types}\ \to \ \textrm{Output Types})[\textrm{Parameter Types}]
\end{align}

For the particular case of the simple weighted linear layer, we would express this signature as

\begin{align}(\mathbb{R}\times\mathbb{R}\to\mathbb{R})[\mathbb{R}, \mathbb{R}]
\end{align}

\section{Standard Library}

Physika features a standard library of mathematical functions that should be familiar from other programming environments

\begin{lstlisting}[language=python]
abs, max, min, uniform, zeros
\end{lstlisting}

\end{document}